\section*{Chapitre \ref{ch:g7:sediments-and-fossils} --- Sédiments et fossiles}

\begin{solution}{exo:g7:sediments-and-fossils:1}
(a) \omterm{def:g7:sediments-and-fossils:fossil}{moule}~; (b) \omterm{def:g7:sediments-and-fossils:fossil}{trace fossile}~; (c) \omterm{def:g7:sediments-and-fossils:sediment}{sédiment}~; (d) \omterm{def:g7:sediments-and-fossils:fossil}{moulage}~; (e) \omterm{def:g7:sediments-and-fossils:layer}{strate}~;
(f) \omterm{def:g7:sediments-and-fossils:fossil}{fossile}.
\end{solution}

\begin{solution}{exo:g7:sediments-and-fossils:2}
(b) l'animal meurt~; (d) la \omterm{def:g3:stones-pebbles-sand:sand}{boue} recouvre vite la coquille~; (c) de
nouvelles couches s'entassent et la coquille se change en pierre~; (a)
l'\omterm{def:g5:water-shapes-the-land:erosion}{érosion} use les couches qui la recouvrent.
\end{solution}

\begin{solution}{exo:g7:sediments-and-fossils:3}
(a) Une \omterm{def:g2:water-on-earth:ocean}{mer} chaude, peu profonde et claire. (b) Une eau peu profonde, au
bord d'un \omterm{def:g1:landscapes-around-us:water}{lac}, d'une \omterm{def:g1:landscapes-around-us:water}{rivière} ou d'une \omterm{def:g2:water-on-earth:ocean}{mer}. (c) Une forêt marécageuse ou un
marais. (d) La terre ferme, couverte de plantes (et survolée par des
libellules), près d'une eau où la \omterm{def:g3:stones-pebbles-sand:sand}{boue} s'est déposée.
\end{solution}

\begin{solution}{exo:g7:sediments-and-fossils:4}
(a) Faux~: la plupart des êtres vivants ne laissent aucun \omterm{def:g7:sediments-and-fossils:fossil}{fossile}~; ils
pourrissent ou se dissolvent avant d'être enfouis. (b) Vrai~: c'est une
\omterm{def:g7:sediments-and-fossils:fossil}{trace fossile}. (c) Faux~: les grains les plus gros et les plus lourds se
déposent en premier, la \omterm{def:g3:stones-pebbles-sand:sand}{boue} fine en dernier. (d) Vrai.
\end{solution}

\begin{solution}{exo:g7:sediments-and-fossils:5}
(a) Grès \qty{120}{m}, calcaire \qty{180}{m}, argilite \qty{100}{m}.
(b) $120 + 180 + 100 = \qty{400}{m}$. (c) $180 \div 400 = 0.45$~:
\textbf{45~\%}. (d) Son sommet a été usé par l'\omterm{def:g5:water-shapes-the-land:erosion}{érosion}~: ce qui reste n'est
pas la couche entière, et ce qui en reste change d'un endroit à l'autre.
\end{solution}

\begin{solution}{exo:g7:sediments-and-fossils:6}
(a) De bas en haut~: le gravier, puis le \omterm{def:g3:stones-pebbles-sand:sand}{sable}, puis la \omterm{def:g3:stones-pebbles-sand:sand}{boue} au-dessus,
sous l'eau claire. (b) Les minuscules grains de \omterm{def:g3:stones-pebbles-sand:sand}{boue} tombent très
lentement~: ils sont encore en train de tomber au bout d'une heure. (c) Une
seconde série de couches au-dessus de la première~: du \omterm{def:g3:stones-pebbles-sand:sand}{sable}, puis de
nouveau de la \omterm{def:g3:stones-pebbles-sand:sand}{boue}. (d) Le gravier, tout au fond.
\end{solution}

\begin{solution}{exo:g7:sediments-and-fossils:7}
Le plus de chances~: (b), une coquille dure, vite enfouie sous la \omterm{def:g3:stones-pebbles-sand:sand}{boue}
fine d'une eau calme. Ensuite (d)~: la coquille est dure et enfouie, mais
une \omterm{def:g5:water-shapes-the-land:flood}{crue} peut aussi la déplacer et la briser. Le moins de chances~: (a),
un corps mou sans parties dures, sur des rochers où aucune \omterm{def:g3:stones-pebbles-sand:sand}{boue} ne se
dépose pour le recouvrir. (c) ne vaut guère mieux~: l'os reste à découvert
et il est décomposé et usé.
\end{solution}

\begin{solution}{exo:g7:sediments-and-fossils:8}
(a) Un creux dont les parois portent la forme de la coquille~: un \omterm{def:g7:sediments-and-fossils:fossil}{moule}.
(b) Un \omterm{def:g7:sediments-and-fossils:fossil}{moulage}, une copie en pierre de la coquille. (c) Non~: la coquille
s'est dissoute~; le \omterm{def:g7:sediments-and-fossils:fossil}{moulage} est fait de nouveaux \omterm{def:g6:rocks-and-minerals:mineral}{minéraux}. (d) L'eau passe
facilement entre les gros grains de \omterm{def:g3:stones-pebbles-sand:sand}{sable}~: elle peut donc atteindre la
coquille et la dissoudre~; dans la \omterm{def:g3:stones-pebbles-sand:sand}{boue} fine, elle passe beaucoup moins
facilement.
\end{solution}

\begin{solution}{exo:g7:sediments-and-fossils:9}
(a) $1.6 \times 1000 = \qty{1600}{m}$. (b) $1600 \div 80 = 20$ couches.
(c) $160 \div 1600 = 0.1$~: \textbf{10~\%}.
\end{solution}

\begin{solution}{exo:g7:sediments-and-fossils:10}
(a) Non~: un pied ne laisse d'empreinte que dans un sol mou, du \omterm{def:g3:stones-pebbles-sand:sand}{sable}
mouillé ou de la \omterm{def:g3:stones-pebbles-sand:sand}{boue}~; le sol est devenu de la \omterm{def:g3:stones-pebbles-sand:rock}{roche} plus tard. (b) La
vague ou la \omterm{def:g5:water-shapes-the-land:flood}{crue} suivante les efface. (c) Un nouveau \omterm{def:g7:sediments-and-fossils:sediment}{sédiment} devait les
recouvrir et les remplir rapidement, avant qu'elles ne soient effacées.
(d) Une \omterm{def:g7:sediments-and-fossils:fossil}{trace fossile}~: elle montre ce que l'animal a fait, et non une
partie de son corps.
\end{solution}

\begin{solution}{exo:g7:sediments-and-fossils:11}
(a) $20 + 5 = \qty{25}{mm}$. (b) $1000 \div 25 = 40$ paires~: la carotte
enregistre \textbf{40 \omterm{def:g5:water-shapes-the-land:flood}{crues}}. (c) Le \omterm{def:g3:stones-pebbles-sand:sand}{sable} est lourd~: il se dépose peu
après l'entrée de la \omterm{def:g1:landscapes-around-us:water}{rivière} dans le \omterm{def:g1:landscapes-around-us:water}{lac}, près de l'\omterm{def:g5:water-shapes-the-land:source-mouth}{embouchure}~; seule une
petite partie va loin. (d) Au large, il ne reste dans l'eau que la \omterm{def:g3:stones-pebbles-sand:sand}{boue}
fine~: elle s'y dépose lentement, quand l'eau est calme.
\end{solution}

\begin{solution}{exo:g7:sediments-and-fossils:12}
(a) Les coquilles sont dures~: elles durent assez longtemps pour être
enfouies, et beaucoup d'animaux à coquille vivaient sur le fond de la \omterm{def:g2:water-on-earth:ocean}{mer},
là où la \omterm{def:g3:stones-pebbles-sand:sand}{boue} se dépose sur eux. (b) Les corps mous pourrissent avant
d'avoir pu être enfouis. (c) Non~: les animaux mous ne laissent presque
jamais de \omterm{def:g7:sediments-and-fossils:fossil}{fossiles}~; la liste tirée des \omterm{def:g7:sediments-and-fossils:fossil}{fossiles} est donc incomplète~:
elle montre surtout les animaux à coquille ou à os.
\end{solution}

\begin{solution}{pb:g7:sediments-and-fossils:1}
\textbf{1.} Le grès, le calcaire, l'argilite, la craie (la première en
bas de la colonne).

\textbf{2.} \qty{120}{m}, \qty{180}{m} et \qty{100}{m}~; au total
\qty{400}{m}.

\textbf{3.} $180 \div 400 = 0.45$~: \textbf{45~\%}.

\textbf{4.} Son sommet a été usé~: la partie haute de la couche a disparu,
et personne ne peut mesurer ce qui n'existe plus.

\textbf{5.} En eau peu profonde, au bord d'un \omterm{def:g1:landscapes-around-us:water}{lac}, d'une \omterm{def:g1:landscapes-around-us:water}{rivière} ou d'une
\omterm{def:g2:water-on-earth:ocean}{mer}~: les rides de courant se forment sur le \omterm{def:g3:stones-pebbles-sand:sand}{sable} en eau peu profonde, et
les empreintes se laissent dans du \omterm{def:g3:stones-pebbles-sand:sand}{sable} ou de la \omterm{def:g3:stones-pebbles-sand:sand}{boue} mous et humides,
souvent sur un rivage.

\textbf{6.} Dans une \omterm{def:g2:water-on-earth:ocean}{mer} chaude, peu profonde et claire~: les coraux.

\textbf{7.} Calme~: la \omterm{def:g3:stones-pebbles-sand:sand}{boue} fine ne se dépose que là où l'eau est presque immobile.

\textbf{8.} D'abord, le lieu était un rivage, avec une eau peu profonde et
un \omterm{def:g3:stones-pebbles-sand:sand}{sable} mou sur lequel marchaient des animaux. Puis la \omterm{def:g2:water-on-earth:ocean}{mer} l'a recouvert~:
une \omterm{def:g2:water-on-earth:ocean}{mer} chaude et peu profonde, où poussaient des coraux et où
s'entassaient des coquilles. Plus tard, l'eau est devenue plus calme,
peut-être plus profonde, et seule de la \omterm{def:g3:stones-pebbles-sand:sand}{boue} fine s'est déposée, avec les
coquilles des ammonites, des animaux parents des calmars et des pieuvres.

\textbf{9.} $180 \times 100 = 18\,000$~cm.

\textbf{10.} $120~\text{m} = 12\,000$~cm, et $12\,000 \times 20 =
240\,000$~ans.

\textbf{11.} $1~\text{m} = 100$~cm, et $100 \times 100 = 10\,000$~ans.

\textbf{12.} $180 \times 10\,000 = 1\,800\,000$~: environ
\textbf{1\,800\,000 ans}, presque deux millions d'années. C'est
$1\,800\,000 \div 240\,000 = 7.5$ fois plus long que pour le grès.
\end{solution}
